% Einheit 3 von 4 – Ebenen im Koordinatensystem (bank/ebenen/e3.jsonl, zusammenbau v0.1)
%% TODO Zeitmarke Einheit 3: Marken-Zeile nicht lesbar
%% TODO Zweigzeile Teil 1: was hier gelernt wird (Satz in Schülersprache aus dem Einheitstitel) – Einheit 3
\einheitenkopf[e3]{Einheit 3 von 4 · Ebenen im Koordinatensystem}
\zweigzeile{Abitur GK}
\setcounter{aufgabe}{14}

%% TODO Titel: Ich-kann-Satz fehlt in der Bank (Platzhalter: Kettenname) – Nr. 15
%% TODO Anweisung: ein Satz über der ersten Teilaufgabe fehlt in der Bank – Nr. 15
\begin{aufgabe}{Besondere Lagen}
%% TODO \rechenplatz{n}: Schrittzahl des Grundfalls fehlt in der Bank (nur bei fester Schreibform, 2.3 b)
\begin{teile}
\teil $E\colon 4x + 5z = 30$. Welche Variable fehlt? \\ \kreuz{$x$} \\ \kreuz{$y$} \\ \kreuz{$z$} \\ \kreuz{keine}
\teil Gleichung der $xz$-Ebene?
\teil Begründe, dass $E\colon 2y + 3z = 18$ parallel zur $x$-Achse verläuft. \hfill (Abitur 2020 LK)
\teil Welche besondere Lage hat $E\colon 4x - 5y = 0$? Beschreibe sie. \hfill (Abitur 2022 GK)
\teil Berechne die Spurpunkte (Schnittpunkte mit den Achsen) von $E\colon 3x + y + 2z = 6$ und zeichne $E$ als Spurdreieck ins Koordinatensystem.

\begin{ksys3}[xyz]\end{ksys3}
\teil Ein Ball fliegt auf der Bahn $X_t(18 - 6t | 4 | -5t^2 + 8t + 1)$; $X_t$ ist sein Ort nach $t$ Sekunden (1 LE = 1 m). Gib die besondere Lage der Ebene $L$ an, in der die Bahn verläuft, und eine Gleichung von $L$. \hfill (Abitur 2023 LK)
\teil Ein Körper ist symmetrisch zur $yz$-Ebene. Ein Weg verläuft in der Ebene $L\colon 6x - 5y = 0$, die senkrecht auf der $xy$-Ebene steht. Verläuft der Weg in der Symmetrieebene? \hfill (Abitur 2026 LK)
\teil Der Punkt $P(4 | 0 | 3)$ ist Eckpunkt eines Rechtecks. Die Gerade $g\colon \vec x = (2 | 5 | 1) + t \cdot (0 | 3 | 0)$, $t \in \mathbb{R}$, verläuft orthogonal zur Ebene des Rechtecks. Begründe, dass das Rechteck in der $xz$-Ebene liegt. \hfill (Abitur 2018 LK)
\teil Eine Kletterwand ist ein ebenes Viereck mit $A(2 | 6 | 1)$, $B(6 | 4 | 1)$, $C(5 | 4{,}5 | 5)$ und $D(1 | 6{,}5 | 5)$; die $xy$-Ebene ist der Boden (1 LE = 1 m). Weise nach, dass die Wand vertikal steht. \hfill (Abitur 2022 GK)
\teil Ein Körper ist symmetrisch zur $xz$-Ebene. Die Ebene seiner Seitenfläche $ABCD$ steht senkrecht auf dieser Symmetrieebene, und $(4 | k | 7)$ ist ein Normalenvektor dieser Ebene. Begründe ohne Rechnung, dass $k = 0$ ist. \hfill (Abitur 2026 LK)
\end{teile}
\end{aufgabe}

%% TODO Titel: Ich-kann-Satz fehlt in der Bank (Platzhalter: Kettenname) – Nr. 16
%% TODO Anweisung: ein Satz über der ersten Teilaufgabe fehlt in der Bank – Nr. 16
\begin{aufgabe}{Besondere Lagen – Fehler finden, Begründen, Anwendung, Darstellungswechsel}
\begin{teile}
\teil Lea sagt: „$E\colon 5x + 2z = 30$ enthält die $y$-Achse, weil $y$ fehlt.“ Finde den Fehler.
\teil Begründe, warum eine Ebene, in deren Koordinatengleichung $z$ nicht vorkommt, parallel zur $z$-Achse ist.
\teil Eine Hallenwand liegt in der Ebene $4x + 3y = 36$ (1 LE = 1 m, die $xy$-Ebene ist der Boden). Steht die Wand senkrecht auf dem Boden, und wo trifft ihre Unterkante die $x$-Achse?
\teil Zeichne ein rechteckiges Stück der Ebene $z = 3$ ins Koordinatensystem.

\begin{ksys3}[xyz]\end{ksys3}
\end{teile}
\end{aufgabe}
